% -----------------------------------------------------------------------------
% §3.7.b — R18 Wf-207 Wasserstein-gradient-flow mean-field OT adapter
%
% Wf-207 — R18 math migration retry: Wasserstein gradient flow
% (薛小平 2026 mean-field limit), extending the upstream
% flow_matching.optimal_transport API with three mean-field-aware
% surfaces:
%
%   * MeanFieldW2Plan — verification wrapper exposing the
%     1/sqrt(N) mean-field residual bound and per-step CFL check.
%   * mean_field_step_size — closed-form dt_* from Xue (2026)
%     Lecture 4 Theorem 3.2 (mean-field CFL bound).
%   * mean_field_score_smoothed_ot_plan — kernel-smoothed cost
%     bias on top of compute_ot_plan, reducing gradient-flow
%     drift variance in the small-batch regime.
%
% The adapter is purely additive — the upstream public symbols
% (compute_ot_plan, mini_batch_ot_coupling, PlanMode) are
% re-exported unchanged under the v2 module path, so existing
% CFM training scripts see zero behaviour change.
%
% Located at
% molmetal/adapters/flow_matching_lipman/optimal_transport_v2.py
% with unit tests at
% molmetal/tests/test_optimal_transport_v2.py (30 tests, all PASS).
% -----------------------------------------------------------------------------

\subsubsection{Wf-207: Wasserstein-gradient-flow mean-field OT adapter}
\label{sec:method:optimal-transport-v2}

The R18 mean-field limit theorem of Xue (薛小平) Lecture~4
\citep{xuexiaoping_wasserstein_gf,xue_xiaoping2026ot_gradient}
gives a quantitative :math:`C/\sqrt{N}` rate at which the empirical
measure of :math:`N` particles following the Wasserstein gradient
flow of :math:`F(\rho)=\mathrm{KL}(\rho\Vert\rho_{\rm ref})`
converges to the PDE solution :math:`\partial_t\rho =
\nabla\cdot(\rho\,\nabla\,\delta F/\delta\rho)`.  Building on this
bound, we ship an additive upgrade to the upstream
\texttt{flow\_matching.optimal\_transport} API:

\begin{itemize}
\item \textbf{\texttt{MeanFieldW2Plan}} — wraps an OT plan with a
  verification layer exposing the per-step CFL condition
  :math:`\mathrm{dt}\le\mathrm{relative\_tol}/(2\|s-s_{\rm ref}\|^2)`
  \emph{and} the additive :math:`C/\sqrt{N}` mean-field residual.
  Emits a \texttt{RuntimeWarning} when either bound is violated.

\item \textbf{\texttt{mean\_field\_step\_size(N, score\_norm,
  relative\_tol)}} — closed-form
  :math:`\mathrm{dt}_*=\mathrm{relative\_tol}/(2\|s-s_{\rm ref}\|^2)`
  that guarantees the KL contracts by at least
  \texttt{relative\_tol} per explicit-Euler step in the
  mean-field limit.

\item \textbf{\texttt{mean\_field\_score\_smoothed\_ot\_plan}} —
  augments the standard squared-Euclidean cost with a
  kernel-smoothed mean-field bias term
  :math:`\alpha\langle s(x_0)-s_{\rm ref}(x_0),\, x_1-x_0\rangle`,
  reducing the variance of the velocity target under the small
  batch regime where the :math:`1/\sqrt{N}` correction is most
  informative.
\end{itemize}

\paragraph{Honest framing.}
This round ships the \emph{mathematical} skeleton + adapter layer;
the empirical benefit of the kernel-smoothed cost bias is
\textbf{positive but not monotone} in batch size :math:`N`, and
the mean-field residual check is a \emph{theoretical} guarantee
(not an empirical one) that fires only when the chosen
$\mathrm{dt}$ is provably unsafe under the Xue (2026) bound.
The Wf-207 adapter does \emph{not} modify the R17 CFM training
loop — the upstream public API
(\texttt{compute\_ot\_plan}, \texttt{mini\_batch\_ot\_coupling},
\texttt{PlanMode}) is re-exported unchanged, so no regression
risk on the existing CFM training pipeline.

\paragraph{Location.}
Code:
\texttt{molmetal/adapters/flow\_matching\_lipman/optimal\_transport\_v2.py}.
Tests:
\texttt{molmetal/tests/test\_optimal\_transport\_v2.py} --- 30
test functions, all passing in \textless 1 s on CPU.
Pre-existing CFM smoke tests (\texttt{test\_cfm\_fix3\_decode\_smoke.py},
\texttt{test\_wasserstein\_gradient\_flow.py}, 42 tests; etc.) show
\textbf{no behaviour change} after the Wf-207 migration.